%%%PAGE 2 rot=0 legibility=good summary=原子物理章节提纲(思维导图式)：光电效应与波粒二象性、原子结构与原子核、玻尔理论与能级跃迁、核衰变、核反应与核能等专题清单
%%%BLOCK id=002-01 ch=17 sec=光电效应·波粒二象性 kind=summary alt=none cont=none y=0.03-0.27
\textbf{原子物理}
\begin{itemize}
\item 光电效应\quad $E_k=h\nu-W_0=\dfrac{hc}{\lambda}-W_0$\quad $h\nu_c=W_0$\quad $eU_c=E_k$\quad $E_k=\dfrac12m_ev^2$
\item 爱因斯坦光电效应方程\quad 光子 $\varepsilon=h\nu$\quad $h=6.63\times10^{-34}\unit{J\cdot s}$\quad $hc=2\times10^{-25}$ % CHECK: 首字形似“受”，按语义转为“爱”；hc 后原稿未写单位
\item 光的波粒二象性
  \begin{itemize}
  \item[$\to$] 干涉、衍射、偏振 % 原稿：该箭头从“光的波”上方引出
  \item[$\to$] 光电效应、康普顿效应、光有动量、能量、 % 原稿：该箭头从“粒”下方引出
  \end{itemize}
\end{itemize}
%%%END
%%%BLOCK id=002-02 ch=17 sec=光电效应·规律 kind=knowledge alt=none cont=none y=0.27-0.345
\textbf{光电效应，波粒二象性} % 页面左侧标注(总标题，对应下面三个专题：光电效应规律、光电效应图像、波粒二象性与物质波)

光电效应规律的理解及应用.
\begin{itemize}
\item[] \unclear{入射}光强度大 $\to$ 光子数目多 $\to$ 发射光电子多 $\to$ 饱和电流大
\item[] 光子频率高 $\to$ 光子能量大 $\to$ 光电子的最大初动能大
\end{itemize}
%%%END
%%%BLOCK id=002-03 ch=17 sec=光电效应·图像 kind=diagram alt=none cont=none y=0.345-0.43
光电效应图像的理解及应用.

%FIG p002_a 0.228 0.3695 0.382 0.4225
\notefig[0.25]{figures/p002_a.png}
\noindent $k=h$\qquad $W_0=|-E|$

%FIG p002_b 0.508 0.3705 0.642 0.4258
\notefig[0.25]{figures/p002_b.png}

% 下图为I--U曲线(黄光、蓝光两条)；图左上手写的“h=ke”已含在此图中
%FIG p002_c 0.65 0.352 0.87 0.4205
\notefig[0.3]{figures/p002_c.png}
\noindent 高强光 左大频 % CHECK: “高强光”首字形似“高”，语义上疑为“同”；另I--U图中“黄光”曲线(上)起点比“蓝光”更靠左，与“左大频”似有矛盾，原样转录
%%%END
%%%BLOCK id=002-04 ch=17 sec=波粒二象性·物质波 kind=knowledge alt=none cont=none y=0.42-0.50
\textbf{对波粒二象性、物质波的理解.}
\begin{itemize}
\item 数量上\quad 少 粒子性\quad 多 波动性
\item 频率上\quad 高 粒子性\quad 低 波动性
\item 传播与作用上\\
  $\uparrow$ 波动性\quad $\uparrow$ 粒子性 % 原稿：两个向上箭头分别在“波动性”“粒子性”上方，指向上一行的“传播与作用上”
\item 波动性和粒子性的统一\quad $\varepsilon=h\nu$\quad $p=\dfrac{h}{\lambda}$
\item 物质波\quad $\lambda=\dfrac{h}{p}=\dfrac{h}{mv}$
\end{itemize}
%%%END
%%%BLOCK id=002-05 ch=17 sec=原子结构·原子核提纲 kind=summary alt=none cont=none y=0.48-0.61
\textbf{原子结构，原子结构} % 页面左侧标注(其右为大括号，含下面“原子结构”“原子核”两支)；CHECK: 第二个“原子结构”或应为“原子核”，原样转录
\begin{itemize}
\item 原子结构\\
  \begin{tabular}[t]{@{}lll@{}}
  汤姆孙发现电子 & 氢原子光谱 & \\
  原子的核式结构 & 玻尔理论 & 氢原子能级、能级公式、 \\
  \end{tabular}
\item 原子核\\
  \begin{tabular}[t]{@{}ll@{}}
  贝克勒尔、天然放射现象、 & 放射性同位素的应用与防护 \\
  原子核的组成 & 核力和核能 \\
  原子核的衰变 & 裂变反应和聚变反应 \\
  \end{tabular}
\end{itemize}
%%%END
%%%BLOCK id=002-06 ch=17 sec=原子的核式结构 kind=note alt=03,04,06,09 cont=none y=0.61-0.67
\textbf{原子的核式结构}

库仑定律、牛顿第二定律、圆周运动公式、功能关系、能量守恒定律、
%%%END
%%%BLOCK id=002-07 ch=17 sec=玻尔理论·能级跃迁 kind=knowledge alt=none cont=none y=0.68-0.87
\textbf{玻尔理论的理解与应用}

\textbf{两类能级跃迁.} % 页面左侧标注，其右为大括号，含下面两项
\begin{itemize}
\item 自发跃迁（高$\to$低）
\item 受激跃迁（低$\to$高）\ $\xrightarrow{\ \text{吸}E\ }$
  \begin{itemize}
  \item 光照（吸收光子）：必须恰好 $h\nu=E_m-E_n$
  \item 碰撞加热.\quad $E_{\text{外}}\ge E_m-E_n$
  \item 吸收大于电离能的光子：原子发生电离 $E_\infty$.
  \end{itemize}
\end{itemize}

% 页面右侧另有一个大括号，含下面三项
\begin{itemize}
\item 电离态
\item $n=\infty$\quad $E=0$. % 原稿“E”前有一个极小的“2”形笔画，疑为逗号/污迹，未转录
\item 电离能\quad 原子从基态或某一激发态跃迁到电离态所需要吸收的最小能量.
\end{itemize}
%%%END
%%%BLOCK id=002-08 ch=17 sec=原子核衰变·半衰期 kind=summary alt=none cont=none y=0.735-0.775
\textbf{原子核的衰变及半衰期}
%%%END
%%%BLOCK id=002-09 ch=17 sec=核反应方程·核能 kind=summary alt=none cont=none y=0.80-0.84
\textbf{核反应方程与核能的计算.}
%%%END
%%%PAGEEND 2
